\section{The original arithmetic specialization theorem}
\label{sec:arithmetic-specialization-v42}

\status{This chapter assembles the exact original
\texttt{SpecIntIff.spec} endpoint and all eleven fields of
$H.\mathrm{IsSpec}$.  It keeps the given algebraic isomorphism,
arithmetic point, algebraic representation, coefficient maps,
finite-part characters, normalized Hecke operators, and one
specialization map across every level.  The necessary coefficient
point argument, elementary torus identities, and level-transport
statements are written below.  No density of arithmetic points,
nonvanishing of a Hecke block, or same-ideal Hida construction
is inferred from this specialization theorem.}

Fix the source data $L,\mathcal\Omega,\ell,n,c,H$, with
$L$ a number field, $\mathcal\Omega/L$ the specified algebraic
closure, $\ell$ prime, $c$ the specified coefficient datum,
and $H:\mathrm{HSet}(L,\mathcal\Omega,\ell,n,c)$ satisfying
all its original $\mathrm{Ax}$ fields.  In particular the
source assumes $\ell\ne2$ and $n\geq1$; these hypotheses
are retained.  Write $O_0=c.O\subset\overline{\mathbb Q}_\ell$.
Let $A$ be an $O_0$-algebra of class $\mathrm{IsCO}$:
it is complete Noetherian local, the structural map is local,
and every residue class of $A$ is represented by an element
of $O_0$.  Fix the actual character $\theta$ with
$H.\mathrm{IsWt}(\theta)$, the specified algebraic field
isomorphism
\[
             \iota:\overline{\mathbb Q}_\ell\simeq\mathbb C,
\]
an $O_0$-algebra map $x:A\to\overline{\mathbb Q}_\ell$,
and an original dominant weight function $\lambda$ such that
$H.\mathrm{Arith}(\theta,x,\lambda)$.
Thus some integer $e_0\geq0$ satisfies
\begin{equation}\label{eq:asp-arithmetic}
 \operatorname{finPart}_{v,i}(y)=1
 \quad\text{whenever}\quad y\in U_c(F_v,e_0),
 \qquad F_v=L_{H.S_l(v)},
\end{equation}
where $y$ is a local unit and the literal finite part is
\begin{equation}\label{eq:asp-finpart}
 \operatorname{finPart}_{v,i}(y)=
       \frac{x\bigl(\theta(e_{v,i}(y))\bigr)}
       {\operatorname{algPow}(H.S_l(v),
                                   \operatorname{lowW}(\lambda,i),y)}.
\end{equation}
Here $e_{v,i}(y)$ is the original diagonal element supported
at the single local coordinate $(v,i)$.  The subgroup
$U_c(F_v,r)$ consists of local units congruent to one
modulo $\mathfrak m_{F_v}^r$.  It is the source $\mathrm{Uc}$,
so the original $T_c(e)$ uses $\mathrm{Uc}(e+1)$.

The subgroup at an auxiliary level $K=(K_q)_{q\in H.Q}$
is the original join
\begin{equation}\label{eq:asp-level}
       U_K=H.\mathrm{lev}(K)
          =H.UQ\vee\bigvee_{q\in H.Q}
                            H.\mathrm{sec}(H.w_q(q))(K_q).
\end{equation}
Every $K_q$ is compact open.  All functions, representations,
and operators in this chapter use the actual Frame $H.X$.

\subsection{Capturing the given arithmetic point in one actual coefficient ring}

% BEGIN ASP1
\paragraph{ASP1. Integral elements and units under the specified field isomorphism.}
\label{claim:ASP1}
For every original coefficient datum $b$ and every
$\mu\in\mathbb C$, with the $b.O$-algebra map given by
the restriction of $\iota$,
\begin{align}
 \mu\text{ is integral over }b.O
       &\quad\Longleftrightarrow\quad|\iota^{-1}(\mu)|_\ell\leq1,
       \label{eq:asp-integral}\\
 \mathrm{IsIntUnit}(b.O,\mu)
       &\quad\Longleftrightarrow\quad|\iota^{-1}(\mu)|_\ell=1.
       \label{eq:asp-unit}
\end{align}
\begin{proof}
Every element of $b.O$ has norm at most one: it is integral
over $\mathbb Z_\ell$ by finite generation, and in a monic
equation a root of norm greater than one would give a
uniquely largest term.  The ultrametric inequality forbids
their sum from being zero.  The same largest-term argument
therefore shows that every root of a monic polynomial with
coefficients in $b.O$ has norm at most one.

Conversely, for $z\in\overline{\mathbb Q}_\ell$ with
$|z|_\ell\leq1$, the field $\mathbb Q_\ell(z)$ is finite
over $\mathbb Q_\ell$.  The norm-unit-ball proof of WRP2
shows that its integer ring is a finite
$\mathbb Z_\ell$-module.  Since it contains $z$, the
determinant trick makes $z$ integral over
$\mathbb Z_\ell$, hence over $b.O$ using the same monic
polynomial.  Transporting the polynomial by $\iota$
proves \eqref{eq:asp-integral} for the specified algebra
map.  In particular this proves the norm-integrality
test in the algebraic closure itself, without replacing
it by a completion.

The original $\mathrm{IsIntUnit}$ means
$\mu\ne0$ and integrality of both $\mu$ and $\mu^{-1}$.
For a nonzero $z$, the two norm inequalities are
$|z|_\ell\leq1$ and $|z|_\ell^{-1}\leq1$, equivalent
to $|z|_\ell=1$.  Norm one already excludes zero.
This proves \eqref{eq:asp-unit}.
\end{proof}
% END ASP1

% BEGIN ASP2
\paragraph{ASP2. The given algebra point has finite coefficients and uniform ideal-power decay.}
\label{claim:ASP2}
There is an original coefficient datum $c_1\geq c$
containing $x(A)$, and there is a real number
$0<\eta<1$ such that
\begin{equation}\label{eq:asp-point-bound}
            |x(a)|_\ell\leq\eta^k
                  \quad(a\in\mathfrak m_A^k, k\geq0).
\end{equation}
In particular $x(\mathfrak m_A)$ consists of elements
of norm strictly less than one.
\begin{proof}
First take $a\in\mathfrak m_A$.  If $|x(a)|_\ell\geq1$,
then $x(a)\ne0$ and its inverse is integral over
$\mathbb Z_\ell$ by ASP1.  Multiply a monic equation
for $x(a)^{-1}$ by the necessary power of $x(a)$ to obtain
\[
             1+b_1x(a)+\cdots+b_rx(a)^r=0,
                                  \qquad b_i\in\mathbb Z_\ell.
\]
The corresponding element $1+b_1a+\cdots+b_ra^r$ of
$A$ lies in $1+\mathfrak m_A$, so it is a unit.  A
unital ring map cannot send a unit to zero, a
contradiction.  Thus $|x(a)|_\ell<1$.

For arbitrary $a\in A$, residue surjectivity supplies
$o\in O_0$ such that $a-o\in\mathfrak m_A$.  Since
$x$ fixes the specified image of $O_0$, and those
elements have norm at most one by ASP1, the
ultrametric inequality gives $|x(a)|_\ell\leq1$.
Choose generators $t_1,\ldots,t_r$ of the Noetherian
ideal $\mathfrak m_A$, and choose
$0<\eta<1$ exceeding their finitely many image norms.
Every element of $\mathfrak m_A^k$ is a finite sum
of products of $k$ generators with coefficients in
$A$.  The bound just proved on those coefficients
gives \eqref{eq:asp-point-bound}.

Let $E\subset\overline{\mathbb Q}_\ell$ be the
compositum of $\operatorname{Frac}(O_0)$ and the fields
generated by $x(t_1),\ldots,x(t_r)$.  It is finite
over $\mathbb Q_\ell$, using GA2 for the first field.
For each $a\in A$ and $N\geq1$ there is a polynomial
$P_N\in O_0[T_1,\ldots,T_r]$ such that
$a-P_N(t_1,\ldots,t_r)\in\mathfrak m_A^N$.
Indeed start with a residue representative for $N=1$.
Given such a polynomial, express its error as a finite
sum of monomials of degree $N$ in the $t_i$, with
coefficients in $A$.  Replace each coefficient by
its representative in $O_0$ modulo $\mathfrak m_A$.
Adding the resulting polynomial improves the error
to $\mathfrak m_A^{N+1}$.

All $x(P_N(t))$ belong to $E$, and by
\eqref{eq:asp-point-bound} they converge to $x(a)$.
The finite extension $E$ is closed in the normed
algebraic closure by the coordinate compactness
argument of GA2.  Thus $x(a)\in E$, and its norm is
at most one.  The actual ring
$O_E=\{z\in E:|z|_\ell\leq1\}$ is a finite
$\mathbb Z_\ell$ DVR by the explicit proof in WRP2.
It contains $O_0$ and $x(A)$; take $c_1.O=O_E$.
This captures the specified map itself, not an
abstract map into an isomorphic coefficient field.
\end{proof}
% END ASP2

% BEGIN ASP3
\paragraph{ASP3. The original weight representation and a local coefficient map at the point.}
\label{claim:ASP3}
There are $c'\geq c_1\geq c$ and an original
$W:\mathrm{WtRep}(H.X.I,c'.O,\mathbb C)$ such that
its matrix representation $\xi$ is explicitly
intertwined with the unchanged $\operatorname{algRep}$
for $(\iota,\lambda)$ and its character satisfies
the original $\operatorname{lowW}$ formula.
Writing $O=c'.O$, the actual coefficient map
\[
                \alpha:A\longrightarrow O,
                        \qquad\alpha(a)=x(a)
\]
is local, and $O\hookrightarrow\mathbb C$ is exactly
$\beta=\iota|_O$.  Both $A$ and $O$ satisfy
the original CNL hypotheses of SMP.
\begin{proof}
Apply WRP12 with the original
$L,\ell,n,H.V,H.S_l,\iota,\lambda$ and the actual
$c_1$ from ASP2.  It constructs the coefficient
extension and all ten weight-representation fields,
with the explicit intertwiner and exact character.
The ring $O$ in that construction is the norm unit
ball of a finite extension, so its maximal ideal
consists precisely of elements of norm less than
one.  ASP2 gives
$\alpha(\mathfrak m_A)\subset\mathfrak m_O$.
Equivalently, if $\alpha(a)$ is a unit, then $a$
is a unit; thus $\alpha$ is the original required
local homomorphism.  Its composite with $\beta$
is exactly $a\mapsto\iota(x(a))$.

WRP2 proves that $O$ is complete Noetherian local
with finite residue field.  The ring $A$ is complete
Noetherian local by $\mathrm{IsCO}$.  The induced
map on residue fields from $O_0$ is injective,
as every unital map of fields is, and surjective
by the given residue representatives.  The
residue field of $O_0$ is finite: $\ell$ is a
nonunit by its norm, so this field is a quotient
of the finite ring $O_0/\ell O_0$.
Hence $A$ has finite residue field as well.
This verifies every coefficient hypothesis
needed to apply SMP to the actual point.
\end{proof}
% END ASP3

\subsection{The finite character and the unchanged classical spaces}

% BEGIN ASP4
\paragraph{ASP4. The actual finite-part torus character and its uniform depth.}
\label{claim:ASP4}
There is a character $\psi:T\to O^\times$ with
$\beta(\psi(t))=\chi(t)$ on the compact torus.
Define the original finite character
\begin{equation}\label{eq:asp-phi}
          \phi(t)=\alpha(\theta(t))\psi(t)^{-1}
                              \quad(t\in T).
\end{equation}
Then
\begin{align}
 \iota(x(\theta(t)))&=\beta(\phi(t))\chi(t),
                                       \label{eq:asp-balance}\\
 \beta(\phi(e_{v,i}(y)))&=
                             \iota(\operatorname{finPart}_{v,i}(y)).
                                       \label{eq:asp-elementary-phi}
\end{align}
For $c_* =\max(e_0,W.c_0)$, $\phi$ is trivial on
$T_c(c_*)$, and
$\operatorname{finPart}_{v,i}(y)=1$ for every
$y\in U_c(F_v,c_*+1)$.
\begin{proof}
The compact-torus integrality field of $W$ and
its zeroth-row identity give a preimage in $O$
of $\chi(t)=\xi(t)_{00}$.  Injectivity of $\beta$
makes that preimage unique.  The preimage for
$t^{-1}$ multiplies it to one, so it is a unit.
The multiplicativity and identity equations for
$\chi$ then imply the same equations in $O$,
again by injectivity.  This defines $\psi$ as
a character into the actual unit group.  Equation
\eqref{eq:asp-balance} follows from the definition
\eqref{eq:asp-phi}.

WRP7 evaluated at the single-coordinate torus
element says
\[
 \chi(e_{v,i}(y))=
       \iota(\operatorname{algPow}(H.S_l(v),
                              \operatorname{lowW}(\lambda,i),y)).
\]
All other factors in its finite product are one.
The denominator is nonzero, since every local
embedding sends the nonzero unit $y$ to a nonzero
element.  Divide \eqref{eq:asp-balance} by this
identity to obtain \eqref{eq:asp-elementary-phi},
with exactly the source finite part.

Every compact diagonal tuple $t$ is the product
over $v,i$ of $e_{v,i}(y_{v,i})$ for its actual
diagonal entries.  By OST26--OST27,
$t\in T_c(c_*)$ means
$y_{v,i}\in U_c(F_v,c_*+1)$ for every entry.
Since $c_*+1\geq e_0$, these units belong to
$U_c(F_v,e_0)$, by inclusion of ideal powers.
Equation \eqref{eq:asp-arithmetic} and
\eqref{eq:asp-elementary-phi} make $\phi$ one
on each factor, by injectivity of $\beta$.
Its product is therefore one.  The same inclusion
of unit congruence groups proves the final
finite-part assertion.  This is the precise
source depth shift by one, uniform in every
$v,i,y$.
\end{proof}
% END ASP4

% BEGIN ASP5
\paragraph{ASP5. Exact equality of the two classical-space definitions.}
\label{claim:ASP5}
For every original auxiliary level $K$ and every
depth $e$,
\begin{equation}\label{eq:asp-classical}
 H.\mathrm{ClAt}(\theta,\iota,x,\lambda,\xi,K,e)
    =H.X.\mathrm{Cl}(H.\tau_C(\iota),\xi,U_K,e,\beta\phi).
\end{equation}
The away representation on the right is exactly
the coefficientwise specialization of $H.\tau_A(A)$.
\begin{proof}
The source $\mathrm{Big}$ condition is the
twisted rational-equivariance condition together
with invariance under $H.UQ$.  Since
$H.UQ\subset U_K$, the separate $U_K$-invariance
condition already supplies this last requirement.
Thus the rational and away conditions on the two
sides are the same.  The upper-unipotent and
depth-$e$ lower-unipotent conditions are literally
the same as well.

The remaining condition on the left prescribes
the eigenvalue
$\iota(\operatorname{finPart}_{v,i}(y))^{-1}$
for each elementary compact torus element
$e_{v,i}(y)$.  By ASP4 this is
$\beta(\phi(e_{v,i}(y)))^{-1}$.
Every element of the compact torus is the finite
product of its elementary diagonal coordinates,
as in ASP4.  Applying the representation to this
product proves that the elementary eigenvalue
conditions imply the eigenvalue condition for
every $t\in T$ with character $(\beta\phi)^{-1}$.
The converse is restriction to those elements.
This proves equality of the actual subspaces.

For the coefficient identity, every entry of
$H.\tau$ belongs to $O_0$.  The map $x$ fixes
the specified $O_0$-structure, so applying
$\iota x$ after $O_0\to A$ is just applying
$\iota$ to that same entry.  This gives
$H.\tau_C(\iota)=(\operatorname{GL}(\iota x))
\circ H.\tau_A(A)$ entrywise, with the original
section maps and away action unchanged.
\end{proof}
% END ASP5

% BEGIN ASP6
\paragraph{ASP6. Exact positive normalizations and equality of ordinary spectra.}
\label{claim:ASP6}
Let $a_{v,j}$ denote the original $\mathrm{alV}(v,j)$,
with the first $j$ local diagonal entries equal to
the original uniformizer and all others equal to
one.  For every $j\geq0$,
\begin{equation}\label{eq:asp-lowval}
                  \chi(a_{v,j})=H.\mathrm{lowVal}(\iota,\lambda,v,j).
\end{equation}
The original local operator indexed by
$(v,j)\in H.V\times\mathrm{Fin}(n)$ is exactly
$T_{a_{v,j+1}}$.  The original
$H.\mathrm{uzC}(\iota,\lambda,W.d+1)$ is exactly $T_z$.
Consequently, for every $K,e$,
\begin{equation}\label{eq:asp-ordinary}
 H.\mathrm{OrdAt}(\theta,\iota,x,\lambda,\xi,K,e)
  =\mathrm{ClOrd}(W,H.\tau_C(\iota),U_K,e,\beta\phi).
\end{equation}
\begin{proof}
In WRP7's original character formula, evaluation
at $a_{v,j}$ leaves precisely the factors
$\operatorname{algPow}(H.S_l(v),
\operatorname{lowW}(\lambda,i),\varpi_v)$
with $i<j$; every other diagonal entry is one.
This is the literal definition of $H.\mathrm{lowVal}$,
giving \eqref{eq:asp-lowval}.  By definition the
source local operator uses the index $j+1$ and
normalizes its original $U$-operator by this value
inverted.  This is exactly $T_{a_{v,j+1}}$.
Each $a_{v,j}$ belongs to the original positive
monoid: its local exponent sequence is one for
$i<j$ and zero afterward, hence is decreasing;
all other components have exponent zero.  This
also covers $j\geq n$, when that local sequence
is constant one.

The diagonal identity for the distinguished element is
\begin{equation}\label{eq:asp-z-factor}
                    z=\prod_{v\in H.V}\prod_{j=0}^{n-1}a_{v,j}.
\end{equation}
Indeed, at coordinate $i$ exactly $n-1-i$ indices
$j\in\{0,\ldots,n-1\}$ satisfy $i<j$, giving
the original exponent $n-1-i$.  All factors are
diagonal and commute.  Apply the character to
\eqref{eq:asp-z-factor} and use
\eqref{eq:asp-lowval}.  The scalar normalizing
$H.\mathrm{uzC}$ is thus exactly $\chi(z)^{-1}$,
with its source product indexed from zero to
$n-1$.  This proves equality of the actual
distinguished operators.

The two definitions in \eqref{eq:asp-ordinary}
sum intersections with maximal generalized
eigenspaces of that same operator.  Their
classical spaces agree by ASP5.  The allowed
eigenvalues on the left satisfy
$|\iota^{-1}(\mu)|_\ell=1$; on the right they
satisfy $\mathrm{IsIntUnit}(O,\mu)$.
These conditions are identical by ASP1.
The sums and all their summands therefore
agree as actual subspaces, proving the claimed
ordinary-space equality.
\end{proof}
% END ASP6

\subsection{The original level family and its exact finite Hecke sums}

% BEGIN ASP7
\paragraph{ASP7. Common refinements and transport by the original auxiliary sections.}
\label{claim:ASP7}
The original family $U_K$ is directed toward smaller
subgroups.  Moreover for every $K$, $q\in H.Q$, and
$g\in\operatorname{GL}_n(L_{H.w_q(q)})$, there is
one original level $K^g$ such that, for every
representation $\rho$ of $H.X.U_b$ over any
coefficient ring,
\begin{equation}\label{eq:asp-level-transport}
 \rho\bigl(H.\mathrm{sec}(H.w_q(q))(g)\bigr)
                     (M^{U_K})\subset M^{U_{K^g}}.
\end{equation}
The choice of $K^g$ depends only on $K,q,g$, not
on the coefficients or the representation.
\begin{proof}
For $K,K'$, take the componentwise intersections
$K''_q=K_q\cap K'_q$.  They are open, and are
compact since an open subgroup is closed and
their intersection is a closed subset of a compact
subgroup.  Thus $K''$ is an original allowed
level.  Monotonicity of images and joins in
\eqref{eq:asp-level} gives
$U_{K''}\subset U_K\cap U_{K'}$.

For transport, replace the $q$ component by
$gK_qg^{-1}$ and leave the other components
unchanged.  Conjugation is a homeomorphism,
so this is again an allowed compact open level.
Put $s=H.\mathrm{sec}(H.w_q(q))(g)$ and let
$v\in M^{U_K}$.  Every element of $H.UQ$
fixes $\rho(s)v$, because it commutes with $s$
by the exact $\mathrm{Q\_UQ}$ field.  At a
component $q'\ne q$, the same assertion follows
from $\mathrm{Q\_Q}$ and invariance of $v$
under the original $K_{q'}$.  At component $q$,
write $h=gkg^{-1}$ with $k\in K_q$.  The
identity $hg=gk$ implies
\[
 \rho(H.\mathrm{sec}(H.w_q(q))(h))\rho(s)v
   =\rho(s)\rho(H.\mathrm{sec}(H.w_q(q))(k))v
   =\rho(s)v.
\]
The stabilizer of a vector under a representation
is a subgroup.  Since it contains every displayed
generator, it contains their exact join $U_{K^g}$.
This proves \eqref{eq:asp-level-transport} for every
representation at once, with the original order
of multiplication.
\end{proof}
% END ASP7

% BEGIN ASP8
\paragraph{ASP8. Original spherical operators are the same finite sums over every coefficient ring.}
\label{claim:ASP8}
For $w\in H.W_p$ and $j\geq0$, there is one
finite indexed family $(u_r)_{r\in R}\subset H.X.U_b$
such that, for every representation $\rho$ of
$H.X.U_b$,
\begin{equation}\label{eq:asp-spherical-sum}
 \mathrm{sphOp}(\rho\circ H.\mathrm{sec}(w),j)
                           =\sum_{r\in R}\rho(u_r).
\end{equation}
It is the original integral double-coset sum, with
no averaging factor.  This sum preserves
$M^{U_K}$ for every original $K$ and every $\rho$.
The same family works for both the $A$-valued
and the complex classical away actions.
\begin{proof}
Put $G_w=\operatorname{GL}_n(L_w)$ and
$C_w=\operatorname{GL}_n(\mathcal O_{L_w})$.
The original orbit $C_wa_{w,j}C_w/C_w$ is finite:
it is the image of the compact group $C_w$ in
the discrete quotient $G_w/C_w$, as proved in
OHA1.  Take this finite set as $R$, use the
original quotient representatives appearing in
the definition of $\mathrm{dc}$, and set
$u_r=H.\mathrm{sec}(w)(r)$.  Unfolding that
definition gives \eqref{eq:asp-spherical-sum}.
Its representatives and multiplicities are fixed
before $\rho$, its coefficient ring, or $K$
is selected.

For preservation, a $U_K$-fixed vector is
$C_w$-fixed because $\mathrm{W\_le}$ puts
$H.\mathrm{sec}(w)(C_w)$ inside $H.UQ\subset U_K$.
OHA1's finite coset permutation argument makes
its spherical image $C_w$-fixed.  Given
$u\in H.UQ$, the original $\mathrm{W\_dec}$
field writes
\[
 u=H.\mathrm{sec}(w)(k)u',\quad k\in C_w,
 \quad u'\in H.UQ,\quad
       [u',H.\mathrm{sec}(w)(g')]=1\ \text{for every }g'.
\]
The action of $u'$ commutes term by term with
the finite sum and fixes the input, so it fixes
the image.  The factor at $k$ fixes the image
by its $C_w$ invariance.  Each auxiliary component
also fixes the image, because the source
$\mathrm{sectionsCommuteAtSplit}$ makes its
action commute with every summand and it fixes
the input.  The stabilizer subgroup therefore
contains the join $U_K$.  This is the same
argument as OHA3, and uses only a group
representation, so it applies equally to both
required coefficient settings and to all their
ambient fixed spaces.
\end{proof}
% END ASP8

\subsection{Installation of the one arithmetic specialization map}

% BEGIN ASP9
\paragraph{ASP9. A single map for the given point and every original level.}
\label{claim:ASP9}
There is one map
\begin{equation}\label{eq:asp-map}
 \mathrm{sp}:H.X.\mathrm{Fn}(A,H.m)\longrightarrow
              \bigl(\mathrm{Fin}(H.m)\to
                              H.X.\mathrm{Fn}(\mathbb C,W.d+1)\bigr)
\end{equation}
with all the conclusions of SMP10 for the actual
$H$-data, the actual point $x$, the fixed $\iota$,
and the depth $c_*$ from ASP4.
\begin{proof}
Apply SMP10 with the original Frame $H.X$,
source ring $A$, target ring $O=c'.O$, and the
specified maps
$\alpha(a)=x(a)$ and $\beta=\iota|_O$ from ASP3.
Their composite is precisely $\iota x$.
The weight representation is the actual $W$
of ASP3, and $\tau=H.\tau_A(A)$.
The character $\theta$ has the required
continuity and free-stabilizer fields by the
given $H.\mathrm{IsWt}(\theta)$.
ASP4 supplies $\phi$, the character balance,
and its triviality on $T_c(c_*)$, with
$c_*\geq W.c_0$.

The family of subgroups is exactly
$K\mapsto H.\mathrm{lev}(K)$.  It is directed
by ASP7, is finite-level by $H.\mathrm{Ax}.\mathrm{fin}$,
and has the free condition at each $K$ by
$H.\mathrm{IsWt}.\mathrm{free}$.
Finally ASP5 identifies the complex away
representation with the required coefficient
specialization.  Thus every SMP hypothesis
has been verified on the original objects.
Take its map defined by the unique good lift.
It is defined on the whole ambient function
space and is fixed once for this arithmetic
point and weight.  In particular it is not
selected anew at a level, depth, operator,
basis, or ideal precision.
\end{proof}
% END ASP9

% BEGIN ASP10
\paragraph{ASP10. Classical finiteness, stability, and the original norm bounds.}
\label{claim:ASP10}
For every original $K$ and $e\geq c_*$, the
space $H.\mathrm{ClAt}(\theta,\iota,x,\lambda,\xi,K,e)$
is finite dimensional.  Every original operator
$H.\mathrm{opC}(\iota,\lambda,W.d+1,i)$ preserves it.
For every nonzero vector $F$ in that space and
every eigenvalue equation
\[
             H.\mathrm{opC}(i)F=\mu F,
\]
one has the original bound $|\iota^{-1}(\mu)|_\ell\leq1$.
\begin{proof}
ASP5 identifies the space with the classical
space in SMP.  Its finite dimensionality is
SMP2.  At a local index $(v,j)$ the operator
is $T_{a_{v,j+1}}$ by ASP6, and SMP2 supplies
preservation.  At an away index $(w,j)$,
ASP8 writes it as the finite sum for
$a_{w,j+1}$.  That sum preserves the ambient
$U_K$-fixed space by ASP8.  Away translations
also preserve twisted left equivariance and
commute with all local right translations,
so their sum preserves the other classical
conditions; this is the explicit argument of
SMP9.  Hence it preserves the same original
classical space.

For a nonzero eigenvector, SMP3 proves
integrality over $O$ for either its normalized
positive operator or its finite away sum.
The latter conclusion did not require
preservation of the original space, although
that property has just been checked here.
ASP1 converts the integral assertion under
the same fixed $\iota$ to exactly the required
norm bound.  This proves the original
$\mathrm{fin}$, $\mathrm{stab}$, and
$\mathrm{int}$ fields with every $K,e,i,F,\mu$
quantifier retained.
\end{proof}
% END ASP10

% BEGIN ASP11
\paragraph{ASP11. All algebraic and Hecke properties of the actual map.}
\label{claim:ASP11}
For the single map \eqref{eq:asp-map}, the
original $\mathrm{add}$, $\mathrm{smul}$,
$\mathrm{span}$, $\mathrm{indep}$, and
$\mathrm{op}$ fields of $H.\mathrm{IsSpec}$ hold.
Explicitly it is additive on every original
ordinary module and satisfies
\[
                \mathrm{sp}(bf)=\iota(x(b))\mathrm{sp}(f).
\]
Its complex span at every $K,e\geq c_*$ is
the literal $H.\mathrm{OrdAt}(K,e)$, every
original $A$-basis has complex-linearly-independent
images, and for every original Hecke index $i$,
\begin{equation}\label{eq:asp-hecke-intertwining}
              \mathrm{sp}(H.\mathrm{op}(A,i)f)
                  =H.\mathrm{opC}(\iota,\lambda,W.d+1,i)
                                                   \mathrm{sp}(f)
\end{equation}
whenever $f$ belongs to the original ordinary module
at $K$.
\begin{proof}
SMP6 gives additivity and semilinearity with scalar
$\beta(\alpha(b))=\iota(x(b))$.
SMP7 gives independence of every original
$A$-basis and the exact classical ordinary span
at every level and admissible depth.  ASP6
identifies that ordinary space with the literal
$H.\mathrm{OrdAt}$, including its original
norm-one eigenvalue condition.

For a local index $(v,j)$, the two operators
in \eqref{eq:asp-hecke-intertwining} are the
original $U_{a_{v,j+1},A}$ and
$T_{a_{v,j+1}}$ by definition and ASP6.
SMP8 therefore gives the equality.  For an
away index $(w,j)$, use ASP8's one fixed
finite list at $w,j+1$ for both coefficient
settings.  Its complex sum satisfies the
full fixed-space preservation hypothesis of
SMP9 with $K'=K$, by ASP8.  Applying SMP9
gives exactly the two sides of
\eqref{eq:asp-hecke-intertwining}, with the
original integral spherical normalization.
This verifies every operator index rather
than only the distinguished ordinary operator.
\end{proof}
% END ASP11

% BEGIN ASP12
\paragraph{ASP12. The original auxiliary-group action clause.}
\label{claim:ASP12}
For all original $K,K'$, $q\in H.Q$, and
$g\in\operatorname{GL}_n(L_{H.w_q(q)})$,
if $f$ is ordinary at $K$ and its original
away translate by $s=H.\mathrm{sec}(H.w_q(q))(g)$
is ordinary at $K'$, then
\begin{equation}\label{eq:asp-action}
       \mathrm{sp}\bigl(H.X.\operatorname{ra}_{H.\tau_A(A)}(s)f\bigr)
          =H.\operatorname{atW}(\iota,W.d+1,H.w_q(q))(g)
                                                       \mathrm{sp}(f).
\end{equation}
\begin{proof}
Take the original conjugated level $K^g$ of ASP7.
Applied to the complex ambient away representation,
\eqref{eq:asp-level-transport} says that the singleton
sum consisting of $s$ carries the whole
$U_K$-fixed space into the $U_{K^g}$-fixed space.
SMP9 therefore gives the required equality for
this singleton sum on every ordinary $f$ at $K$.
By its actual definition the complex away action
at $s$ is
$H.\operatorname{atW}(\iota,W.d+1,H.w_q(q))(g)$,
giving \eqref{eq:asp-action}.
The equality is an equality of ambient functions
for the one map already constructed.  Thus it
also holds with the additional source premise
that the $A$-valued translate is ordinary at the
specified $K'$, whether or not $K'=K^g$.
No level-dependent choice of specialization is
made in using that premise.
\end{proof}
% END ASP12

% BEGIN ASP13
\paragraph{ASP13. The full original \texttt{SpecIntIff.spec} endpoint.}
\label{claim:ASP13}
For every original input in the setup there
exist $D\in\mathbb N$, a matrix representation
$\xi:\mathrm{GV}(n,H.S_l)\to\operatorname{GL}_D(\mathbb C)$,
a depth $c_*$, and the one map $\mathrm{sp}$
on the original ambient functions such that
\[
          H.\mathrm{IsSpec}(\theta,\iota,x,\lambda,
                                                \xi,c_*,\mathrm{sp})
\]
holds.  This is exactly
$\mathrm{UpperTheory.V.Hi.spec}$ in the pinned
$\mathrm{SpecIntIff.lean}$.
\begin{proof}
Take $D=W.d+1$ and $\xi=W.\xi$ from ASP3,
$c_*=\max(e_0,W.c_0)$ from ASP4, and the
single map of ASP9.  All eleven assertion
fields of the original $H.\mathrm{IsSpec}$
are furnished individually as follows:
\begin{center}
\begin{tabular}{ll}
Original field & Written proof \\
\hline
$\mathrm{alg}$ & ASP3, using the exact WRP intertwiner \\
$\mathrm{triv}$ & ASP4, with the original $c_*+1$ \\
$\mathrm{fin}$ & ASP10 \\
$\mathrm{stab}$ & ASP10 \\
$\mathrm{int}$ & ASP10 and the fixed-$\iota$ norm test \\
$\mathrm{add}$ & ASP11 \\
$\mathrm{smul}$ & ASP11, scalar $\iota(x(b))$ \\
$\mathrm{span}$ & ASP11 and ASP6 \\
$\mathrm{indep}$ & ASP11, every original $A$-basis \\
$\mathrm{op}$ & ASP11, every spherical and local index \\
$\mathrm{act}$ & ASP12, every original auxiliary section
\end{tabular}
\end{center}
The algebraic representation, finite character,
complex coefficient map, local uniformizers,
ordinary eigenvalue condition and Hecke
normalizations were identified explicitly in
ASP3--ASP6.  The coefficient-point capture
and its local map were proved for the
specified $x$, not merely for some arithmetic
point with the same residue.  The family of
levels and each finite spherical sum were
matched before applying the single-map
construction.  These observations verify the
literal consumer contract at all its
quantifiers and finish the original endpoint.
\end{proof}
% END ASP13

\subsection{Completed endpoint and remaining global work}

The original generic arithmetic specialization
theorem is now supplied in full human-readable
form, conditional only on its printed original
$H.\mathrm{Ax}$, $H.\mathrm{IsWt}$ and
$H.\mathrm{Arith}$ inputs.  For the actual
reconstruction datum, FNL16 supplies the
original $H.\mathrm{Ax}$, and the earlier
weight-character construction supplies
$H.\mathrm{IsWt}$ where invoked.  This
chapter preserves the individual arithmetic
point and original coefficient ring throughout;
the extension $c\leq c_1\leq c'$ and the
specialized map $A\to c'.O$ are explicit.

Specialization alone does not assert that
there are enough arithmetic points to separate
a Hecke algebra, that a chosen localized
ordinary block is nonzero at an arithmetic
point, or that a Hida construction uses the
same original ideal with its full automorphy
and Galois conditions.  Those are distinct
remaining statements.  No such statement,
and no whole RH or GRH obligation, is retired
by the endpoint just proved.
